% META: {"id": "1NSI-ARCHOS-EVAL-A", "chapitre": "1NSI-ARCHITECTURE-OS", "type_objet": "evaluation", "capacites": ["P-ARCH-01A", "P-ARCH-01B", "P-ARCH-03A", "P-ARCH-03C"], "mode_creation": "ex_nihilo", "sources_inspiration": [], "duree_min": 45, "status": "verified"}
\section*{Évaluation A --- Architecture et systèmes d'exploitation}
\textit{Durée : 45 minutes. Ordinateur avec interpréteur Python autorisé.}

\baremeIndicatif{Ex. 1 : 10 pts (von Neumann, instructions) ; Ex. 2 : 10 pts (OS, permissions)}

\bigskip

\textbf{Exercice 1} --- \textit{Von Neumann et instructions machine} \hfill (10 points)

\begin{enumerate}
  \item (3 pts) Nommer les quatre constituants principaux de l'architecture de von Neumann, et préciser lequel stocke à la fois instructions et données.
  \item (4 pts) On ajoute au simulateur \lstinline{executer} du cours l'instruction \lstinline{("MUL", reg_dest, reg_src)} (multiplication : \lstinline{reg_dest <- reg_dest * reg_src}). Écrire le code de cette instruction.
  \item (3 pts) Dérouler l'exécution de ce programme et donner l'état final de \lstinline{memoire} :
\begin{python}
memoire = [6, 7, 0]
programme = [
    ("LOAD", "R0", 0),
    ("LOAD", "R1", 1),
    ("MUL", "R0", "R1"),
    ("STORE", "R0", 2),
    ("HALT",),
]
\end{python}
\end{enumerate}

% BEGIN-VERIFY
% def executer(programme, memoire):
%     registres = {}
%     for instruction in programme:
%         op = instruction[0]
%         if op == "LOAD":
%             _, reg, adresse = instruction
%             registres[reg] = memoire[adresse]
%         elif op == "MUL":
%             _, reg_dest, reg_src = instruction
%             registres[reg_dest] = registres[reg_dest] * registres[reg_src]
%         elif op == "STORE":
%             _, reg, adresse = instruction
%             memoire[adresse] = registres[reg]
%         elif op == "HALT":
%             break
%     return memoire, registres
% memoire = [6, 7, 0]
% programme = [
%     ("LOAD", "R0", 0),
%     ("LOAD", "R1", 1),
%     ("MUL", "R0", "R1"),
%     ("STORE", "R0", 2),
%     ("HALT",),
% ]
% memoire_finale, registres = executer(programme, memoire)
% assert memoire_finale == [6, 7, 42]
% END-VERIFY

\bigskip

\textbf{Exercice 2} --- \textit{Système d'exploitation et permissions} \hfill (10 points)

\begin{enumerate}
  \item (3 pts) Citer deux fonctions assurées par un système d'exploitation.
  \item (4 pts) Un fichier a les permissions \lstinline{"rwxr-x---"}. En utilisant \lstinline{droit_accorde} du cours, déterminer si le propriétaire peut l'exécuter, si le groupe peut l'exécuter, et si les autres peuvent le lire.
  \item (3 pts) Convertir ces permissions en notation octale avec \lstinline{permissions_vers_octal}.
\end{enumerate}

% BEGIN-VERIFY
% def droit_accorde(permissions, categorie, action):
%     index = {"proprietaire": 0, "groupe": 3, "autres": 6}
%     decalage = {"lecture": 0, "ecriture": 1, "execution": 2}
%     position = index[categorie] + decalage[action]
%     return permissions[position] != "-"
% def permissions_vers_octal(permissions):
%     groupes = [permissions[0:3], permissions[3:6], permissions[6:9]]
%     valeurs = []
%     for g in groupes:
%         v = (4 if g[0] == "r" else 0) + (2 if g[1] == "w" else 0) + (1 if g[2] == "x" else 0)
%         valeurs.append(v)
%     return "".join(str(v) for v in valeurs)
% perm = "rwxr-x---"
% assert droit_accorde(perm, "proprietaire", "execution") is True
% assert droit_accorde(perm, "groupe", "execution") is True
% assert droit_accorde(perm, "autres", "lecture") is False
% assert permissions_vers_octal(perm) == "750"
% END-VERIFY
